\subsubsection{2D Intermediate Resolved-Flow Models}
\label{subsubsec:2d-models}

A 2D model is treated here as a conceptual intermediate resolved-flow tier,
usually posed as an axisymmetric or planar flow problem. It retains a spatially
resolved velocity and pressure field on an idealized reduced domain rather than
only the section-averaged state $(A,Q,p_{\mathrm g})$. Within those
idealizations, it can be used to inspect radial velocity structure, throat
acceleration, simplified recirculation, and wall-shear trends that a 1D
area--flow model cannot resolve.

The same idealization also fixes the boundary of the evidence. Axisymmetric and
planar models suppress eccentric plaques, curvature, torsion, branch geometry,
non-axisymmetric secondary flow, and fully patient-specific WSS\@. Their wall
outputs depend on the declared stress tensor, wall location, normal, tangent,
and output operator, just as in the shear convention of
Convention~\ref{conv:appendix-shear-diagnostic}. A 2D result is therefore not
interchangeable with matched 3D or resolved-FSI evidence. The broad continuum
and dimension-reduction background for this reading is supplied by
\cite{GaldiEtAl2008HemodynamicalFlows,KopplHelmig2023DimensionReducedModeling}.

The 2D tier is a diagnostic bridge: useful for controlled mechanism checks, but
too idealized to serve as clinical or resolved-3D evidence. In this report it
can help interpret whether a reduced trend is plausibly tied to throat
acceleration or simplified recirculation, but it is not the selected forward
generator and does not validate the 1D map by itself.
